\documentclass[11pt]{article}
\usepackage[T1]{fontenc}
\usepackage{lmodern,amsmath,amssymb,amsthm,geometry,microtype,xcolor}
\geometry{margin=1in}
\usepackage[colorlinks=true,allcolors=blue!45!black]{hyperref}
\newtheorem{theorem}{Theorem}
\title{A four-card check of the printed shelf-shuffle spectrum}
\author{Myeongjin Kang}
\date{October 25, 2026}
\begin{document}
\maketitle
\begin{abstract}
The printed nonzero-spectrum conjecture for a single-shelf shuffle includes
the value $1/16$ when the deck has four cards. We prove that this value is
not an eigenvalue of the four-card position matrix. A Lean certificate
checks both the matrix obtained by enumerating all fair-coin choices and
its agreement with the source's explicit binomial formula, then excludes
$1/16$ over the complex field. This is an index-range correction to the
conjecture as printed, not a claim of a general new spectral formula.
\end{abstract}
\section{The printed claim}
In the single-shelf model of Clay~\cite{Clay}, cards initially labeled
$1,\ldots,n$ are drawn from the bottom of the deck. Each drawn card is
placed on the top or bottom of the current pile, with equal probability.
The position matrix has entries
\[
 M_{ij}=\frac{\binom{i-1}{j-1}+\binom{i-1}{n-j}}{2^i}
 \quad (1\leq i,j\leq n).
\]
Conjecture 2 of the cited version, also MathDB369281~\cite{MathDB}, lists
the nonzero eigenvalues as
\[
 2^{-2i},\qquad 0\leq i\leq\left\lceil\frac{n-1}{2}\right\rceil.
\]
For $n=4$, the stated upper endpoint is $2$, so the list includes $1/16$.

\begin{theorem}
For a four-card deck, $1/16$ is not an eigenvalue of the position matrix.
Consequently the nonzero-spectrum conjecture is false as printed.
\end{theorem}
\begin{proof}
Substitution into the binomial formula gives
\[
 M=\begin{pmatrix}
 1/2&0&0&1/2\\
 1/4&1/4&1/4&1/4\\
 1/8&3/8&3/8&1/8\\
 1/8&3/8&3/8&1/8
 \end{pmatrix}.
\]
Suppose $Mv=v/16$ for a complex column vector $v=(a,b,c,d)^T$.
Multiplying by $16$ gives
\begin{align*}
 7a+8d&=0,\\
 4a+3b+4c+4d&=0,\\
 2a+6b+5c+2d&=0,\\
 2a+6b+6c+d&=0.
\end{align*}
Subtracting the third equation from the fourth gives $c=d$.
Twice the second equation minus the fourth therefore gives
$6a+9d=0$. Together with $7a+8d=0$, this implies $5d=0$,
then $a=b=c=d=0$. There is no nonzero eigenvector, proving the theorem.
\end{proof}
\section{Verification of the shuffle model}
The certificate does not merely postulate the displayed numerical matrix.
It enumerates all $2^4=16$ Boolean choices for top-versus-bottom placement,
processes cards $4,3,2,1$, and counts the final position of each card.
The resulting count matrix is
\[
 \begin{pmatrix}8&0&0&8\\4&4&4&4\\2&6&6&2\\2&6&6&2\end{pmatrix}.
\]
The first placement coin has no effect on an empty pile; retaining it
double-counts every effective outcome equally, which preserves the
uniform distribution. Division by $16$ yields $M$ exactly. Lean checks
this finite enumeration and its agreement with the source formula by
kernel reduction, without a native decision axiom.

The file \texttt{formal/ShuffleSpectrum.lean} then proves that any complex
vector satisfying $Mv=v/16$ is zero. It expresses the missing eigenvalue
using mathlib's \texttt{Module.End.HasEigenvalue}, and explicitly negates
the universal printed spectral claim using $n=4$ and $i=2$.
The proof includes the nonzero-eigenvector condition and no unproved
spectral premise.

\section{Scope}
This note targets the ceiling endpoint in the version cited. It does not
disprove a conjecture with a different endpoint, claim to determine the
complete spectrum for arbitrary deck size, or claim that this correction
was previously unknown. In particular, changing the upper endpoint is a
different statement and requires its own proof. The certificate is a
complete refutation of the printed statement, with a transparent
four-card example. The neighboring quadratic-eigenvector conjecture is a
separate question.

The implementation uses Lean 4.34.1 and the campaign's pinned mathlib
revision. All local declarations are audited for standard logical axioms
only. This is an AI-assisted verification note, not external peer review.
Original prose is CC BY 4.0 and original proof code is MIT.
\begin{thebibliography}{9}
\bibitem{Clay} A.~Clay, \emph{Guessing Strategies for Shelf-Shuffling
Machines}, arXiv:2507.10294v1, July 14, 2025. Theorem 1.1 and Conjecture 2.
\url{https://arxiv.org/html/2507.10294v1}.
\bibitem{MathDB} MathDB, \emph{The conjectured nonzero spectrum of the
position matrix}, problem 369281. Accessed October 25, 2026.
\url{https://mathdb.com/p/369281}.
\end{thebibliography}
\end{document}
